% notes/field-guide.tex: field guide for the data captures (compiled to notes/field-guide.pdf).
\documentclass[11pt,a4paper]{article}
\usepackage[T1]{fontenc}
\usepackage{lmodern}
\usepackage[british]{babel}
\usepackage[a4paper,margin=22mm]{geometry}
\usepackage{booktabs,array,tabularx}
\usepackage{enumitem}
\usepackage{xcolor}
\usepackage{fancyvrb}
\usepackage[hidelinks]{hyperref}
\setlength{\parskip}{4pt}\setlength{\parindent}{0pt}
\setlist{itemsep=1pt,topsep=2pt}
\definecolor{boxgrey}{HTML}{F1EFE8}
\newcommand{\cmd}[1]{\texttt{#1}}
\newcommand{\notebox}[1]{\par\smallskip\noindent\colorbox{boxgrey}{\begin{minipage}{\dimexpr\linewidth-2\fboxsep\relax}#1\end{minipage}}\par\smallskip}
\fvset{fontsize=\small,frame=single,framesep=2mm,xleftmargin=2mm,xrightmargin=2mm}

\title{\textbf{Field Guide: Capturing the Thesis Data}\\[2pt]\large LTE downlink occupancy with the RTL-SDR Blog V4, Benin City}
\author{Version 2, 2 November 2026 (checked against the finished capture software)}
\date{}

\begin{document}
\maketitle
\vspace{-8mm}

\notebox{%
\textbf{What you will collect.} A band survey (once, day and night), then at least \textbf{12 capture sessions} over at least \textbf{4 days}, by day and at night, from at least \textbf{2 positions} in Benin City. A session is 8 captures of 30\,s each and takes about 20--30 minutes. It uses about 0.9\,GB of disk, so 12 sessions need about 11\,GB plus backups. The settings come from \cmd{notes/method-design.md} and the commands from the \cmd{occupancy} package in \cmd{code/}.
}

\section{Equipment}
\begin{tabularx}{\linewidth}{@{} >{\raggedright\arraybackslash}p{4.2cm} X @{}}
\toprule
Item & Notes \\
\midrule
RTL-SDR Blog V4 & The only receiver used in the study. Never swap in another dongle. \\
Antenna for 700--960\,MHz & The same antenna for every session. It is passive, so the bias tee stays \textbf{off}. \\
50\,$\Omega$ SMA terminator & Fitted in place of the antenna for the Class~A captures. \\
Laptop (Dell Latitude 5490) & Charged or on mains power; at least 15\,GB free disk space. \\
USB extension cable & Optional: moves the dongle away from the laptop if the laptop causes interference. \\
External drive & At least 20\,GB, for a backup after every session. \\
Phone or camera, notebook & For the photographs (Section~\ref{sec:photos}) and notes. \\
Thermometer & Optional: temperature is recorded if known. \\
\bottomrule
\end{tabularx}

\section{One-time preparation (before the first field day)}
All commands are run from the \cmd{code/} folder of the project:
\begin{Verbatim}
cd ~/projects/msc-thesis/masters-thesis/code
\end{Verbatim}
\begin{enumerate}
\item \textbf{Check the receiver.} Plug in the V4 and run \cmd{rtl\_test -t}. It should find the device, report the R828D tuner and print ``RTL-SDR Blog V4 Detected'' with the 29 gain steps. The last line, ``No E4000 tuner found, aborting'', is normal: that test is only for an older tuner. Then check that the study's software can see it, and note the gain steps it offers:
\begin{Verbatim}
uv run python -m occupancy capture --info
\end{Verbatim}
\item \textbf{Check the survey tool.} \cmd{rtl\_power -h} should print its help; it comes with the RTL-SDR driver package.
\item \textbf{Check the cell-search tool.} LTE-Cell-Scanner is installed in \cmd{\textasciitilde/tools/LTE-Cell-Scanner} (installation: thesis Appendix~C). Its program \cmd{CellSearch} should print its options:
\begin{Verbatim}
cd ~/tools/LTE-Cell-Scanner/build && ./src/CellSearch -h
\end{Verbatim}
\item \textbf{Rehearse without the dongle.} The rehearsal writes only to \cmd{data\_synthetic/}, which the thesis never uses:
\begin{Verbatim}
uv run python -m occupancy session --synthetic
\end{Verbatim}
\item \textbf{Clock.} Make sure the laptop sets its time automatically. Capture times are recorded in UTC.
\end{enumerate}

\section{Step 1: band survey (07:00 and 22:00)}
Do this at P1, the parlour balcony, before any capture session, once starting at 07:00 (\cmd{--label day}) and once starting at 22:00 (\cmd{--label night}), the same times as the sessions. Fix one gain for the whole survey: 40.2\,dB, the tuner step nearest the proposal's 40\,dB, so that the recorded gain is the gain applied. Each run takes about half an hour of sweeps (three ranges of 10 minutes), plus the cell searches:
\begin{Verbatim}
uv run python -m occupancy survey --range 758M:803M --gain 40.2 \
    --integration 10 --duration 10m --label day --position P1
uv run python -m occupancy survey --range 790M:894M --gain 40.2 \
    --integration 10 --duration 10m --label day --position P1
uv run python -m occupancy survey --range 925M:960M --gain 40.2 \
    --integration 10 --duration 10m --label day --position P1
\end{Verbatim}
At 22:00, repeat with \cmd{--label night}. The CSV files and plots are saved in \cmd{data/survey/}.
\begin{enumerate}
\item \textbf{Find the carriers.} An LTE carrier shows as a flat plateau of raised power, 5--20\,MHz wide. Note each plateau's centre and width.
\item \textbf{Confirm with the cell scanner.} For each strong plateau, search 0.5\,MHz either side of its centre, at a fixed gain that is one of the tuner's steps (e.g.\ 40.2\,dB), with two tries per frequency. For a plateau near 796\,MHz:
\begin{Verbatim}
cd ~/tools/LTE-Cell-Scanner/build
./src/CellSearch -s 795.5e6 -e 796.5e6 -g 40.2 -n 2
\end{Verbatim}
A plateau counts as a carrier only if a cell is found. From the results table record: \textbf{CID} (cell identity, PCI), \textbf{fc} (carrier centre), \textbf{nRB} (bandwidth: 6, 15, 25, 50, 75, 100 = 1.4, 3, 5, 10, 15, 20\,MHz), \textbf{RXPWR} and \textbf{DPX} (should be FDD). Never use \cmd{LTE-Tracker}; it decodes more than the study needs.
\item \textbf{Choose the carrier.} It should be confirmed by the scanner, strong and steady by day and night, and preferably 10\,MHz wide or more. Only carriers of 10\,MHz or more allow the offset captures; the session helper checks this.
\item \textbf{Choose the idle (Class~B) window.} It is a stretch of at least 1.92\,MHz with no carrier, as close to the chosen carrier as possible. The scanner finds a carrier only at its centre, so do not search the window itself: check instead that the window lies outside every nearby carrier (centre $\pm$ half its bandwidth). The session helper checks this against the chosen carrier. The survey analysis suggests only stretches with no signal at all, narrow or intermittent ones included. The thin spikes evenly spaced about 2.2\,MHz apart are made by the dongle at the centre of each tuning step; ignore them. The survey analysis can rank candidates:
\begin{Verbatim}
uv run python -m occupancy survey --parse data/survey/<file>.csv --carrier 796M
\end{Verbatim}
\item \textbf{Choose the two gains, $g_1$ and $g_2$.} They are fixed for the whole study.
  \begin{itemize}
  \item $g_1$ is the highest gain step at which a capture centred on the carrier is not flagged for clipping (more than 0.1\% of samples at the limits).
  \item $g_2$ is lower, about 10\,dB below $g_1$ (a suggestion, not a fixed rule).
  \item Make trial captures under the session ID \cmd{trial}. Afterwards, move \cmd{data/raw/trial} to a folder outside \cmd{data/} (e.g. \cmd{data\_trials/}) so they are never used.
  \end{itemize}
\item \textbf{Send me the survey plots, the scanner output, and your choice of carrier, Class~B window, $g_1$ and $g_2$} before Session~1. They go into Appendix~A of the thesis.
\end{enumerate}

\section{Step 2: capture sessions}
\subsection*{Plan}
\textbf{Positions:} P1 is the parlour balcony and P2 the bedroom balcony (they face different directions). At each one, mount the antenna at a marked spot, at the same height and orientation every time, at least 30\,cm from metal railings or burglary-proof bars, and photograph the setup; the laptop stays indoors with the cable through the door. \textbf{Before the first P2 session} (S03 in the plan), check P2 once: run \cmd{CellSearch} on the chosen carrier there, and one 10-minute survey of the range that contains it with \cmd{--position P2 --label check}; send me the plot.

Use the same carrier, Class~B window, antenna and gains $g_1$, $g_2$ in every session. Day sessions start at \textbf{07:00} and night sessions at \textbf{22:00}; the session helper warns if you start more than two hours away from these times, and such a session is reported separately. There is one day and one night slot per day, so 12 sessions take at least 6 days. A possible plan, balanced over times and positions:

\begin{center}
\begin{tabular}{@{} l l l l @{\hspace{12mm}} l l l l @{}}
\toprule
Session & Day & Start & Position & Session & Day & Start & Position \\
\midrule
S01 & 1 & 07:00 & P1 & S07 & 4 & 07:00 & P2 \\
S02 & 1 & 22:00 & P1 & S08 & 4 & 22:00 & P1 \\
S03 & 2 & 07:00 & P2 & S09 & 5 & 07:00 & P1 \\
S04 & 2 & 22:00 & P2 & S10 & 5 & 22:00 & P2 \\
S05 & 3 & 07:00 & P1 & S11 & 6 & 07:00 & P2 \\
S06 & 3 & 22:00 & P2 & S12 & 6 & 22:00 & P1 \\
\bottomrule
\end{tabular}
\end{center}

\subsection*{Procedure for every session}
\begin{enumerate}
\item \textbf{Set up} at the position: antenna mounted as before at that position (see its photo), dongle plugged in, laptop on. Wait about 5 minutes before the first capture, so the dongle warms up.
\item \textbf{Cell search} on the chosen carrier only (here 796\,MHz), keeping the samples with the session's name. Note the CID; it should be the same cell as in the survey:
\begin{Verbatim}
./src/CellSearch -s 796e6 -g 40.2 -n 2 --recbin S01_cellsearch.bin
\end{Verbatim}
The scanner and the session helper both need the dongle, so run them one after the other.
\item \textbf{Start the guided session} and answer its questions:
\begin{Verbatim}
uv run python -m occupancy session
\end{Verbatim}
\begin{itemize}
\item It asks for: the session ID (S01, S02, \dots); the position; $g_1$ and $g_2$; the frequency correction (enter 0); the antenna; the temperature; notes on weather and surroundings; the cell-search tool; the carrier's centre frequency, PCI and bandwidth; which side of the carrier to use for the offset window (+ or $-$); and the Class~B and Class~A frequencies.
\item It then shows the plan.
\end{itemize}
\item \textbf{Follow the screen for the 8 captures} of 30\,s each, pressing Enter to start each one:
  \begin{enumerate}[label=\alph*.]
  \item Class~A at $g_1$, then at $g_2$: \textbf{terminator fitted}.
  \item Class~B at $g_1$, then $g_2$: \textbf{remove the terminator and connect the antenna}.
  \item Class~C centred at $g_1$, then $g_2$.
  \item Class~C offset at $g_1$, then $g_2$. These are skipped automatically if the carrier is narrower than 10\,MHz.
  \end{enumerate}
\item \textbf{Check after each capture.} The helper prints the clipping fraction and the mean power, and saves a quick-look image.
  \begin{itemize}
  \item If a capture is flagged or something went wrong (a bumped cable, a sudden noise), retake it. The first attempt is kept and marked ``retaken''.
  \item In the quick-looks, Class~A should be flat and the lowest, Class~B flat, and Class~C a raised band.
  \end{itemize}
\item \textbf{Back up} the session folder \cmd{data/raw/S01/} (and so on) to the external drive the same day.
\end{enumerate}

\notebox{%
\textbf{Do not}:
\begin{itemize}
\item change the gain, antenna or position during a session;
\item use automatic gain (the software refuses it);
\item switch the bias tee on;
\item run other heavy programs during a capture;
\item touch the cable while a capture runs.
\end{itemize}
}

\section{Step 3: labelling check (on a later day)}
Run the cell scanner once more on the carrier. Then make quick-look images of every capture and a checklist to fill in:
\begin{Verbatim}
uv run python -m occupancy inspect
\end{Verbatim}
It writes the images and \cmd{labelling\_check.csv} to \cmd{work/inspect/}. For each capture, look at its image and fill in the \cmd{second\_pass\_class} column (A, B, C\_centre or C\_offset), whether you re-ran the cell search, and any notes. Look especially for an idle window that shows a carrier, or an active capture that looks flat. Then report the agreement between the two passes, which goes into the thesis:
\begin{Verbatim}
uv run python -m occupancy inspect --agreement work/inspect/labelling_check.csv
\end{Verbatim}

\section{Photographs and screenshots for the thesis}\label{sec:photos}
Yes, photographs help: they show the examiners the real set-up behind the figures. Please take:
\begin{itemize}
\item the V4 with the antenna connected, as a close-up on a plain background with a coin or ruler for scale;
\item the V4 with the 50\,$\Omega$ terminator fitted (the Class~A set-up);
\item the complete set-up at each position (laptop, dongle and antenna), as a wide shot;
\item how the antenna is mounted at each position, which also helps you repeat it exactly.
\end{itemize}
For screen content, take \textbf{screenshots}, not photos of the screen, e.g. of the real-time demonstration running (I'll tell you when it is ready). The survey plots are produced by the software, so no picture of them is needed. A photo of the laptop's model label is not needed: the thesis lists the hardware, and the software logs the CPU model automatically when it measures speed.

\textbf{Tips:}
\begin{itemize}
\item Landscape, in daylight, without flash glare; keep the original files.
\item No faces, house numbers or anything that reveals a private address.
\item Put them in a \cmd{photos/} folder in the project, with names like \cmd{2026-11-05\_P1\_setup.jpg}.
\end{itemize}

\section{Troubleshooting}
\begin{tabularx}{\linewidth}{@{} >{\raggedright\arraybackslash}p{4.6cm} X @{}}
\toprule
Problem & What to do \\
\midrule
Device not found & Replug the dongle, run \cmd{rtl\_test -t}, close any other SDR program. \\
Capture flagged for clipping & Retake once. If it persists, note it and tell me. Do \emph{not} change $g_1$ or $g_2$ during the study. \\
Messages about lost samples & Close other programs and retake the capture. \\
Cell scanner finds no cell & Do not capture Class~C that day; note it and tell me. \\
Disk nearly full & Move backed-up sessions to the external drive. \\
\bottomrule
\end{tabularx}

\section{Safety and ethics}
Reception is passive: nothing is transmitted. The cell scanner decodes only what every cell broadcasts openly (its identity and bandwidth); no other content is decoded or recorded. Ask permission before using a private place, and take care of your own safety during night sessions.

\section{What to give me after the fieldwork}
\begin{itemize}
\item The \cmd{data/} folder (raw captures with their \cmd{.json} records, quick-looks and the survey). I can work on it directly on the laptop.
\item The cell-scanner outputs, your notes and the photographs.
\end{itemize}
\end{document}
